%% Draft v0.1 — compile on Overleaf. Single self-contained file (no external .bib).
%% Uses the ACM acmart class in `nonacm` mode so it builds without ACM rights metadata.
\documentclass[sigconf,nonacm]{acmart}
\settopmatter{printacmref=false}
\setcopyright{none}
\pagestyle{plain}

\begin{document}

\title{From Summarization to Synthesis: Privacy-Preserving Generation of
Multiple-Aspect Trajectories from Semantic Locations}

\author{[Author list to be completed]}
\affiliation{%
  \institution{ISTI-CNR, Pisa \& collaborators}
  \country{Italy}}
\email{draft@example.org}

\renewcommand{\shortauthors}{Draft v0.1}

\begin{abstract}
Sharing human-mobility data is scientifically valuable but legally and ethically constrained: raw
trajectories are hyper-identifying, and data owners are reluctant to release them. Synthetic data is
the standard escape, yet existing trajectory generators are \emph{geometry-centric}---they reproduce
where people go but discard the rich semantics (activities, points of interest, transport modes) that
make mobility data useful, and the most faithful (deep) generators tend to memorize real traces. We
propose to generate synthetic mobility in a \emph{semantic state space} rather than a raw geometric
one, obtained from \textsc{MAT-Sum}, a location-centric summarization method that abstracts a city
into \emph{semantic locations}. On this substrate we build a compact pipeline: an anonymized semantic
vocabulary with a tunable generalization/anonymity dial, a mechanistic (Markov) sequence generator
with an empirical dwell-time model, an optional deep (LSTM) alternative, and a spatially-coupled
instantiation step that re-materializes trajectories over real geography without copying any real
path. We evaluate on a small private sample (Akkodis, Paris) as a proof-of-concept and, at scale, on
a public dataset of 581 de-identified individuals in Paris. The semantic abstraction yields three
properties that no single existing technique offers together: (i)~\emph{semantic richness}---the
output is a Multiple-Aspect Trajectory, not a coordinate stream; (ii)~\emph{controllable, structural
privacy}---a single interpretable dial spans the fidelity/privacy frontier and the mechanistic
generator produces \textbf{zero verbatim copies} at every setting, with a distance-to-closest-record
at or above the real-to-real baseline; and (iii)~\emph{data efficiency}---fidelity improves
monotonically with the number of individuals while a comparable deep model ``wins'' only by memorizing
at small sample sizes and loses that lead at scale. We argue this makes location-centric semantic
summarization a natural foundation for privacy-preserving mobility-data sharing.
\end{abstract}

\keywords{synthetic mobility data, multiple-aspect trajectories, trajectory privacy, semantic
summarization, data sharing}

\maketitle

\section{Introduction}
Movement data generated by GPS-enabled devices underpins research in urban planning, epidemiology,
transportation and behavioral modeling. Yet the very properties that make such data useful also make
it dangerous to release: individual traces are remarkably unique---a handful of spatio-temporal points
suffices to re-identify most people~\cite{montjoye2013}---so data owners (telecom operators, mobility
companies, municipalities) keep it locked away. This produces a chronic bottleneck: the organizations
holding the data cannot share it, and researchers who could extract value from it cannot obtain it.

\emph{Synthetic data} is the widely-adopted resolution: replace real records with artificial ones
that preserve aggregate structure while breaking the link to real individuals. For trajectories, a
rich literature exists---mechanistic models grounded in statistical laws of human
mobility~\cite{song2010,pappalardo2018,jiang2016}, deep generative
models~\cite{feng2020,rao2020,zhu2023}, and differentially-private synthesizers~\cite{gursoy2018}.
Two gaps persist. First, almost all of these operate on \textbf{geometry}---they generate coordinates
or grid cells, treating semantics (what a place \emph{is}, what activity happens there) as absent or
as a single side attribute. Second, the most faithful generators are prone to \textbf{memorization},
silently reproducing real traces and defeating the privacy goal.

Our starting observation is that a recent summarization method, \textsc{MAT-Sum}~\cite{pugliese2023},
already provides the missing representation. \textsc{MAT-Sum} enriches a geographic area with semantic
aspects (from OpenStreetMap) and merges cells that share a semantic context into \emph{semantic
locations}, turning each trajectory into a weighted sequence over a finite semantic vocabulary. We
make the case that this is not merely a compression device but a \textbf{generative substrate}:
generation performed over semantic locations inherits everything \emph{above the abstraction
line}---semantics, coarse space, coarse time---while resampling everything below it---exact geometry
and timing. This is precisely what buys both \emph{scale} (the model reasons over semantic
\emph{types}, so it can instantiate new, non-copied trajectories) and \emph{privacy} (mapping a raw
point to a generalized semantic location is a form of semantic spatial cloaking).

The motivation is concrete and practical. A data owner such as a mobility company will not hand over
raw data on faith; but a small, non-confidential sample can be used to demonstrate that a semantic,
non-memorizing generator produces useful yet safe output, unlocking controlled access to the full
dataset from which a distributable synthetic version is then generated. This paper takes the first two
steps of that argument: it establishes the method and validates its key properties, first on a small
private sample and then on a public dataset at scale.

\paragraph{Contributions}
\begin{enumerate}
  \item A pipeline that turns \textsc{MAT-Sum}'s semantic locations into a \emph{privacy-preserving
  generative model} of Multiple-Aspect Trajectories, with an interpretable generalization/anonymity
  dial.
  \item A \emph{spatially-coupled} instantiation step that re-materializes semantic sequences over
  real geography, recovering realistic locality without copying real paths.
  \item An empirical study---proof-of-concept on a private sample and validation on 581 public
  individuals---showing the semantic abstraction delivers semantic richness, controllable privacy and
  data efficiency, and that a simple mechanistic generator dominates a deep one on the fidelity/privacy
  trade-off.
\end{enumerate}

\section{Background and Related Work}
\paragraph{Semantic trajectories and summarization}
Semantic (Multiple-Aspect) trajectories enrich raw movement with heterogeneous
aspects~\cite{spaccapietra2008,bogorny2019}. \textsc{MAT-Sum}~\cite{pugliese2023} adopts a
\emph{location-centric} view: it tessellates an area, enriches tiles with OSM aspects (POIs, land use,
public transport), merges tiles with identical semantic context into semantic locations, and
summarizes each trajectory as a weighted sequence over them, quantified by a summarization rate
($S_{\text{rate}}$) and semantic quality (MUITAS~\cite{petry2019}). We reuse \textsc{MAT-Sum} as-is
for map enrichment and repurpose its output for generation.

\paragraph{Mobility generation}
Mechanistic models reproduce empirical laws (visitation frequency, radius of gyration, jump length,
waiting time) via exploration/preferential-return dynamics~\cite{song2010,pappalardo2018}; they are
interpretable and data-efficient. Deep generators~\cite{feng2020,rao2020,zhu2023} achieve higher
geometric fidelity given large corpora, at the risk of memorization. Differentially-private
synthesizers privatize a mobility model (e.g., noisy Markov counts) before sampling~\cite{gursoy2018}.
To our knowledge, generation of \emph{semantically rich} Multiple-Aspect Trajectories---and the use of
a semantic summarization as the generative substrate---is largely unexplored; this is the gap we
target.

\paragraph{Evaluation}
We adopt the fidelity/utility/privacy triad, with utility measured by train-on-synthetic/test-on-real
(TSTR) and privacy by membership inference, distance-to-closest-record (DCR), and copy/uniqueness
analysis.

\section{Research Questions}
\begin{description}
  \item[RQ1 (semantic richness)] Does generating over semantic locations produce Multiple-Aspect
  trajectories whose \emph{semantics} are preserved, in a way geometry-centric generators cannot?
  \item[RQ2 (controllable privacy)] Does the semantic generalization/anonymity dial $(\tau,k)$ provide
  a controllable fidelity$\leftrightarrow$privacy frontier, and does the generator avoid memorizing
  real individuals?
  \item[RQ3 (data efficiency)] How does synthesis quality behave as the number of individuals grows,
  and how does a mechanistic generator compare to a deep one across data regimes---especially the
  small-sample regime typical of restricted-access settings?
  \item[RQ4 (spatial realism \& transfer)] Can semantic sequences be re-materialized into
  spatially-realistic trajectories over real geography, and does the abstraction enable placement
  beyond the training region? \textit{[cross-city transfer: pending NYC run]}
\end{description}

\section{Methodology}

\subsection{Overview: the abstraction line}
Let a dataset of raw trajectories over a geographic area be given. \textsc{MAT-Sum} maps it to (a)~a
set of \emph{semantic locations} and (b)~for each trajectory, an ordered, temporally-weighted sequence
of the semantic locations it visits. Our generator operates on (b) over the vocabulary (a). Everything
\emph{above the abstraction line}---the semantic identity of places and the coarse spatio-temporal
structure---is modeled and preserved; everything below it---exact coordinates and timestamps---is
discarded and later resampled. Preserved properties are therefore exactly those \textsc{MAT-Sum}
retains, which its own $S_{\text{rate}}$/MUITAS analysis quantifies; freed degrees of freedom are what
enable both scaling and privacy.

\subsection{Anonymized semantic vocabulary (Phase 0)}
From the semantic-mapped trajectories we form ordered \emph{visit sequences} by run-length-collapsing
consecutive points that fall in the same semantic context, accumulating their dwell time. Each
semantic location is represented by a \emph{symbol}: the set of its top-$m$ TF-IDF semantic aspects.
The granularity $m$ controls the vocabulary size (few, coarse symbols for small $m$; many, fine
symbols for large $m$). We then enforce \emph{$k$-anonymity on the vocabulary}: any symbol occurring
fewer than $k$ times is merged into the most Jaccard-similar frequent symbol, guaranteeing every
symbol is supported by at least $k$ occurrences and eliminating singleton, identifying contexts. The
pair $(m,k)$ is the method's \textbf{privacy/utility dial}---the analogue of \textsc{MAT-Sum}'s
similarity threshold $\tau$---governing semantic resolution, and hence jointly governing compression,
semantic quality and privacy.

\subsection{Sequence generators (Phase 1)}
\paragraph{Mechanistic (Markov)}
We fit a first-order Markov chain over the symbol alphabet with Laplace smoothing, together with an
empirical per-symbol dwell-time distribution. Trajectory length is drawn from the empirical length
distribution (a DITRAS-style decoupling of \emph{how many} from \emph{where}), which reproduces
realistic lengths without a termination state whose probability would be diluted over a large
alphabet. Transition and dwell counts can be privatized with calibrated noise for a formal
$\varepsilon$-DP guarantee, which we evaluate as a prototype in \S\ref{sec:alt}.

\paragraph{Deep (LSTM)}
As a learned alternative we train a small LSTM language model over the same alphabet and the same
dwell model, isolating the sequence model. This tests whether a higher-capacity generator earns its
fidelity honestly or by memorizing training individuals.

\paragraph{Variable-order and differentially-private variants}
Because a first-order chain is parameterized exactly by bigram counts, it already matches the bigram
distribution near-optimally; to capture longer-range structure we also consider a \emph{variable-order}
model (VOMM) that interpolates orders $1..K$ in a Jelinek--Mercer fashion, which by construction cannot
underperform order~1 on bigrams and improves higher-order statistics. Orthogonally, for a formal
guarantee we consider a \emph{differentially-private} order-1 model: each individual's transition
contribution is clamped to bound the histogram $L_1$ sensitivity, and Laplace noise is added to the
counts, yielding user-level $\varepsilon$-DP. Both variants are evaluated in \S\ref{sec:alt}.

\subsection{Spatial instantiation (Phases 2--3)}
To obtain trajectories over real geography, each generated symbol is realized at a concrete enriched
tile of that semantic type.

\emph{Decoupled (Phase 2).} Independently for each visit, choose a tile of the required symbol with an
exploration/preferential-return gravity weighting. This preserves coarse spatial spread but, because
the semantic sequence is spatially blind, breaks micro-locality: the nearest tile of a given fine
type may be far, inflating jump lengths.

\emph{Spatially-coupled (Phase 3).} We instead generate directly at the tile level with a transition
that couples three factors---learned semantic compatibility $T[s_{\text{cur}},s_{\text{next}}]$,
spatial proximity $\exp(-d/\text{scale})$, and a visitation prior---wrapped in an EPR explore/return
process~\cite{song2010}. Locality then \emph{emerges} from the process rather than being imposed.
Optionally, consecutive locations are connected with road-network routing (OSRM) to produce
road-consistent, plausible-but-invented GPS geometry; because tiles are chosen anew, no real path is
reproduced.

\subsection{Privacy as controllable semantic generalization}
Mapping raw points to generalized semantic locations is a semantic form of the spatial-cloaking /
$k$-anonymity primitive of location privacy~\cite{gruteser2003}, but generalizing by \emph{meaning}
rather than a blind grid, so it preserves utility. The $(m,k)$ dial sets the generalization strength;
generating \emph{new} sequences (rather than resampling real ones) addresses the sequence-level
uniqueness threat; and the spatial-instantiation and routing steps re-introduce realism without
re-introducing the original data.

\subsection{Evaluation metrics}
\emph{Fidelity}: distributional distances (total variation / Jensen--Shannon) between real and
synthetic on semantic unigram/bigram distributions and MUITAS; since an order-1 model saturates the
bigram objective, we also report \emph{trigram}-TV to expose higher-order generators, plus spatial
radius of gyration and jump length. \emph{Utility}: TSTR on a semantic downstream task \textit{[pending]}. \emph{Privacy}:
exact- and near-copy rate (MUITAS $\geq 0.95$) against the training set; distance-to-closest-record
($\text{DCR}=1-$ nearest-neighbor MUITAS) versus the real-to-real baseline; membership inference
\textit{[pending at scale]}.

\section{Experimental Validation}

\subsection{Datasets}
\paragraph{Akkodis / Paris (proof-of-concept, private)}
A small confidential sample of 20 GPS trajectories in Paris (4{,}039 points). Used only to establish
the pipeline; not a basis for quantitative claims.

\paragraph{Paris-581 (validation, public)}
The Paris subset of the semantically-enriched OSM mobility dataset of~\cite{pugliese2025}: raw GPS
traces merged \emph{one trajectory per de-identified OpenStreetMap contributor}, i.e.\ \textbf{581
individuals}, 1.57\,M points. Because merging is per user, the privacy unit is the person and
per-trajectory metrics are per-person. (A spatial home-clustering heuristic gave a fragile 260--478
estimate contradicting this authoritative count and was discarded; the residual multi-account confound
is rare and noted as a limitation.) Timestamps are real (2007--2024); only user ids are de-identified.
We ran \textsc{MAT-Sum} map-enrichment on this area (squares, zoom~17): 32{,}697 enriched tiles
$\rightarrow$ 18{,}554 semantic locations. \textit{[New York (18{,}765 individuals): pending]}

\subsection{Summarization at scale (context for RQ1)}
On Paris-581, \textsc{MAT-Sum} attains a dataset summarization rate $S_{\text{rate}}=0.921$
(per-individual range 0.45--0.996), within the published Geolife range (0.90--0.98). The
proof-of-concept's lower 0.738 was not a method weakness but data sparsity---the PoC trajectories are
short (84--481 points) while the validation trajectories are densely sampled---confirming that dense
semantic trajectories compress into few semantic locations very effectively.

\subsection{Granularity, fidelity and privacy (RQ2)}
Table~\ref{tab:sweep} sweeps the $(m,k)$ dial on Paris-581 ($N=200$ synthetic; nearest-neighbor over a
150-sample). At every setting the mechanistic generator's DCR sits at or above the real-to-real
baseline---it is never closer to real records than real records are to one another---and verbatim
copies are essentially absent. Increasing $k$ both anonymizes and, by densifying the alphabet,
improves fidelity; the vocabulary grows with the city ($m{=}3\rightarrow959$ symbols, versus 71 on the
small PoC), so fine granularity remains sparse even at $n{=}581$, motivating the scaling study.

\begin{table*}[t]
\caption{$(m,k)$ sweep on Paris-581. $\downarrow$ lower is better, $\uparrow$ higher is better.}
\label{tab:sweep}
\small
\begin{tabular}{@{}rrrrrrrrr@{}}
\toprule
$m$ & $k$ & alphabet & unigram TV\,$\downarrow$ & bigram TV\,$\downarrow$ &
MUITAS$\rightarrow$real\,$\uparrow$ & DCR\,$\uparrow$ & DCR base & copies \\
\midrule
1 & 5 & 52   & 0.021 & 0.097 & 0.913 & 0.087 & 0.079 & 2 \\
2 & 5 & 328  & 0.134 & 0.329 & 0.630 & 0.370 & 0.253 & 0 \\
3 & 5 & 959  & 0.342 & 0.716 & 0.402 & 0.598 & 0.350 & 0 \\
2 & 2 & 441  & 0.211 & 0.412 & 0.571 & 0.429 & 0.274 & 1 \\
3 & 2 & 1410 & 0.463 & 0.830 & 0.346 & 0.654 & 0.388 & 0 \\
\bottomrule
\end{tabular}
\end{table*}

\subsection{Data efficiency (RQ3)}
Fixing the vocabulary ($m{=}2,k{=}5$) and varying the number of training individuals $n$ isolates the
effect of scale (Table~\ref{tab:scaling} and Fig.~\ref{fig:scaling}, 3 seeds). Fidelity improves monotonically---bigram-TV falls
from 0.75 at $n{=}20$ to 0.33 at $n{=}581$---while DCR stays above the real-to-real baseline (0.27) at
every $n$ and copies remain $\approx 0$. The generator is thus data-efficient \emph{and} structurally
leak-free; the small-sample limitation of the PoC is real, and scale addresses it.

\begin{table}[t]
\caption{Data-scaling at fixed vocabulary ($m{=}2,k{=}5$).}
\label{tab:scaling}
\small
\begin{tabular}{@{}rrrrr@{}}
\toprule
$n$ individuals & bigram TV\,$\downarrow$ & DCR synth\,$\uparrow$ & DCR base & copies of train \\
\midrule
20  & 0.750 & 0.496 & 0.267 & 0.0 \\
50  & 0.695 & 0.507 & 0.267 & 0.7 \\
100 & 0.640 & 0.503 & 0.267 & 0.0 \\
300 & 0.459 & 0.453 & 0.267 & 0.0 \\
581 & 0.329 & 0.409 & 0.267 & 0.3 \\
\bottomrule
\end{tabular}
\end{table}

\begin{figure}[t]
\centering
\includegraphics[width=\linewidth]{fig_scaling.png}
\caption{Data-scaling at fixed vocabulary ($m{=}2,k{=}5$). Fidelity (bigram-TV, blue) improves
monotonically with the number of individuals; privacy (DCR, orange) stays above the real-to-real
baseline throughout.}
\label{fig:scaling}
\end{figure}

\subsection{Mechanistic versus deep across scale (RQ3, decisive)}
Training a small LSTM and the order-1 Markov on the \emph{same} subsample across $n$
(Table~\ref{tab:deep}, Fig.~\ref{fig:deep}) gives the sharpest result. The LSTM leads on fidelity at small and mid $n$---but
buys it by \textbf{memorizing}: 14.3\% of its output at $n{=}20$ are near-copies of training
individuals, collapsing toward zero as data grows. The Markov never copies. Moreover the fidelity
advantage narrows and \emph{reverses at $n{=}581$}, where the simple, interpretable, DP-amenable
Markov (bigram-TV 0.401) beats the LSTM (0.441) with zero leakage.

\begin{table}[t]
\caption{Deep (LSTM) vs mechanistic (Markov) across $n$ ($m{=}2,k{=}5$, 2 seeds).}
\label{tab:deep}
\small
\begin{tabular}{@{}rrrrr@{}}
\toprule
$n$ & Markov bi-TV\,$\downarrow$ & LSTM bi-TV\,$\downarrow$ & Markov copies\,\% & LSTM copies\,\% \\
\midrule
20  & 0.779 & 0.674 & 0.0 & 14.3 \\
50  & 0.735 & 0.563 & 0.3 & 1.3 \\
100 & 0.675 & 0.521 & 0.3 & 0.7 \\
300 & 0.500 & 0.439 & 0.0 & 0.3 \\
581 & 0.401 & 0.441 & 0.3 & 0.3 \\
\bottomrule
\end{tabular}
\end{table}

For privacy-preserving synthesis the implication is clear: at small $n$---the restricted-access regime
that motivates the work---the deep model's fidelity is an artifact of copying real people; at scale
the mechanistic model is at least as faithful and never leaks. The mechanistic generator therefore
dominates the fidelity/privacy trade-off across the entire data range studied.

\begin{figure*}[t]
\centering
\includegraphics[width=\textwidth]{fig_deep.png}
\caption{Deep (LSTM) vs mechanistic (Markov) across $n$ ($m{=}2,k{=}5$). Left: the LSTM leads on
fidelity at small/mid $n$ but the Markov overtakes it at $n{=}581$. Right: the LSTM's edge is bought by
memorization---$14\%$ near-copies of training at $n{=}20$, collapsing to $\approx 0$ as data grows;
the Markov never copies.}
\label{fig:deep}
\end{figure*}

\subsection{Alternative generators: pushing fidelity and adding formal privacy}
\label{sec:alt}
Two variants (Table~\ref{tab:alt} and Fig.~\ref{fig:alt}, Paris-581, $m{=}2,k{=}5$) probe whether the mechanistic model can be
improved and made formally private. The interpolated \emph{variable-order} model (VOMM) improves
fidelity while preserving privacy: trigram-TV falls from 0.844 to 0.596 (and bigram-TV from 0.322 to
0.251), with DCR still above the real-to-real baseline (0.244) and zero copies---the generator can be
pushed beyond first order without leaking, and the gain is visible precisely on the higher-order
statistic that a first-order model cannot capture. The \emph{differentially-private} order-1 model
traces the formal-privacy/fidelity frontier: at $\varepsilon{\to}\infty$ it recovers the plain Markov,
and tightening $\varepsilon$ raises privacy (DCR) at a fidelity cost that is steep for this naive
user-level mechanism---a direct consequence of the large per-user sensitivity of trajectory data, and
the reason adaptive DP methods~\cite{gursoy2018,privtrace} are needed. Crucially it renders the privacy
guarantee \emph{formal} rather than empirical, with zero copies throughout.

\begin{table}[t]
\caption{Alternative generators on Paris-581 ($m{=}2,k{=}5$). VOMM improves higher-order fidelity;
DP-Markov adds a formal $\varepsilon$-DP guarantee at a fidelity cost. Real-to-real DCR baseline 0.244.}
\label{tab:alt}
\small
\begin{tabular}{@{}lrrrr@{}}
\toprule
model & bigram TV\,$\downarrow$ & trigram TV\,$\downarrow$ & DCR\,$\uparrow$ & copies \\
\midrule
Markov (order 1)                  & 0.322 & 0.844 & 0.371 & 0 \\
VOMM (interp., $\leq 3$)          & 0.251 & 0.596 & 0.349 & 0 \\
DP-Markov $\varepsilon{=}1$       & 0.964 & 1.000 & 0.656 & 0 \\
DP-Markov $\varepsilon{=}20$      & 0.921 & 0.995 & 0.607 & 0 \\
DP-Markov $\varepsilon{=}100$     & 0.760 & 0.962 & 0.517 & 0 \\
DP-Markov $\varepsilon{=}\infty$  & 0.326 & 0.837 & 0.358 & 0 \\
\bottomrule
\end{tabular}
\end{table}

\begin{figure}[t]
\centering
\includegraphics[width=\linewidth]{fig_altmodels.png}
\caption{DP-Markov privacy/fidelity vs $\varepsilon$ (Paris-581, $m{=}2,k{=}5$). Tightening
$\varepsilon$ (left, more private) raises DCR and worsens fidelity; at $\varepsilon{\to}\infty$ the
plain Markov (dotted) is recovered.}
\label{fig:alt}
\end{figure}

\subsection{Spatial realism (RQ4)}
On the proof-of-concept geography, the decoupled instantiation preserved coarse spatial spread but
inflated jump lengths (real median 0.28\,km vs 1.01\,km) and the radius of gyration (2.36 vs
3.28\,km), because a spatially-blind semantic sequence realized at each type's nearest instance cannot
reproduce micro-locality. The \emph{spatially-coupled} generator recovers it: it matches the radius of
gyration (2.44 vs 2.36\,km) and restores the short-jump mode (0.72\,km), with zero exact path copies (Fig.~\ref{fig:spatial});
the residual jump-length gap against continuously-sampled real GPS is a sampling-granularity artifact.
\textit{[coupled generator at Paris-581 scale: pending]}

\begin{figure*}[t]
\centering
\includegraphics[width=\textwidth]{fig_spatial.png}
\caption{Spatially-coupled generator on the proof-of-concept geography. It recovers the real
short-jump mode (left) that the decoupled variant misses, matches the radius of gyration (middle), and
produces localized home-region trajectories (right).}
\label{fig:spatial}
\end{figure*}

\subsection{Limitations and threats to validity}
(i)~The deep baseline is a single small, untuned LSTM; a larger model could improve fidelity at large
$n$---though not undo the small-$n$ memorization, which is the privacy point. (ii)~Validation so far is
one city; New York provides cross-city robustness and larger-$n$ membership inference. (iii)~Beyond empirical
privacy (DCR, copy-rate, MUITAS) we evaluate a \emph{prototype} user-level $\varepsilon$-DP mechanism
(\S\ref{sec:alt}); a tuned, adaptive DP synthesizer (e.g.\ PrivTrace~\cite{privtrace}) that would shift
the frontier left is left to future work. (iv)~Utility is not yet measured with a TSTR
downstream task. (v)~One real individual with multiple OSM accounts would appear as multiple
``individuals''; this rare confound is inherited from the source dataset.

\section{Conclusion}
We argued and gave first evidence that a location-centric semantic summarization, \textsc{MAT-Sum}, is
a natural substrate for privacy-preserving synthesis of Multiple-Aspect Trajectories. Generating over
a $k$-anonymized semantic vocabulary yields synthetic mobility that is semantically rich, controllably
private and data-efficient---a combination no single existing technique offers---and a simple
mechanistic generator dominates a deep one on the fidelity/privacy trade-off---an interpolated
variable-order variant improves higher-order fidelity while a differentially-private variant adds a
formal guarantee. Natural next steps are cross-city validation (New York), a tuned adaptive-DP
synthesizer, a TSTR utility study, and a membership-inference audit at scale.

%% ---------------------------------------------------------------------------
%% Self-contained bibliography (replace with a .bib + ACM-Reference-Format for submission).
\begin{thebibliography}{99}
\bibitem{bogorny2019} R.~dos Santos Mello, V.~Bogorny, et al. MASTER: A Multiple Aspect View on Trajectories. \emph{Transactions in GIS}, 2019.
\bibitem{feng2020} J.~Feng, Z.~Yang, et al. Learning to Simulate Human Mobility (MoveSim). \emph{KDD}, 2020.
\bibitem{gruteser2003} M.~Gruteser and D.~Grunwald. Anonymous Usage of Location-Based Services through Spatial and Temporal Cloaking. \emph{MobiSys}, 2003.
\bibitem{gursoy2018} M.~E.~G\"ursoy, L.~Liu, et al. Utility-Aware Synthesis of Differentially Private and Attack-Resilient Location Traces (AdaTrace). \emph{ACM CCS}, 2018.
\bibitem{jiang2016} S.~Jiang, Y.~Yang, et al. The TimeGeo Modeling Framework. \emph{PNAS}, 2016.
\bibitem{montjoye2013} Y.-A.~de Montjoye, C.~A.~Hidalgo, et al. Unique in the Crowd: The Privacy Bounds of Human Mobility. \emph{Scientific Reports}, 2013.
\bibitem{pappalardo2018} L.~Pappalardo and F.~Simini. Data-driven Generation of Spatio-temporal Routines in Human Mobility (DITRAS). \emph{Data Mining and Knowledge Discovery}, 2018.
\bibitem{petry2019} L.~M.~Petry, C.~A.~Ferrero, et al. Towards Semantic-Aware Multiple-Aspect Trajectory Similarity Measuring (MUITAS). \emph{IJGIS}, 2019.
\bibitem{pugliese2023} C.~Pugliese, F.~Lettich, F.~Pinelli, C.~Renso. Summarizing Trajectories Using Semantically Enriched Geographical Context (MAT-Sum). \emph{ACM SIGSPATIAL}, 2023.
\bibitem{pugliese2025} C.~Pugliese, F.~Lettich, G.~Rocchietti, C.~Renso, F.~Pinelli. Human Mobility Datasets Enriched With Contextual and Social Dimensions. \emph{arXiv:2510.02333}; Zenodo 15658129, 2025.
\bibitem{privtrace} H.~Wang, Z.~Zhang, et al. PrivTrace: Differentially Private Trajectory Synthesis by Adaptive Markov Models. \emph{USENIX Security}, 2023.
\bibitem{rao2020} J.~Rao, S.~Gao, et al. LSTM-TrajGAN: A Deep Learning Approach to Trajectory Privacy Protection. \emph{GIScience}, 2020.
\bibitem{song2010} C.~Song, T.~Koren, P.~Wang, A.-L.~Barab\'asi. Modelling the Scaling Properties of Human Mobility (EPR). \emph{Nature Physics}, 2010.
\bibitem{spaccapietra2008} S.~Spaccapietra, C.~Parent, et al. A Conceptual View on Trajectories. \emph{Data \& Knowledge Engineering}, 2008.
\bibitem{zhu2023} Y.~Zhu, Y.~Ye, et al. DiffTraj: Generating GPS Trajectory with Diffusion Probabilistic Model. \emph{NeurIPS}, 2023.
\end{thebibliography}

\end{document}
